\documentclass[11pt]{article}
\usepackage[a4paper,margin=1in]{geometry}
\usepackage{graphicx}
\usepackage{booktabs}
\usepackage{float}
\usepackage{subcaption}
\usepackage{hyperref}
\usepackage{amsmath}
\usepackage{array}
\usepackage{longtable}
\hypersetup{colorlinks=true, linkcolor=blue, urlcolor=blue, citecolor=blue}

\title{Efficient Graph Query Processing in PostgreSQL using Recursive SQL and Indexing}
\author{Sameer Patil \and Abhi Jain \and Jatsya Jariwala \and Harsh Suthar}
\date{}

\begin{document}
\maketitle

\section{Introduction}
Graph-structured workloads arise naturally in social networks, recommendation systems, communication graphs, and contact tracing. Although specialized graph databases are available, many real systems continue to rely on relational databases for transactional consistency, operational simplicity, and existing infrastructure. This project studies how far PostgreSQL can support graph query processing when the graph is represented relationally and traversal is implemented using recursive SQL.

The main goal is not to compete with native graph engines on all workloads, but to understand how schema design, recursive common table expressions (CTEs), and indexing affect execution behavior inside a mature relational query processor.

\section{Problem Statement and Objectives}
We model a directed graph using normalized relational tables and implement graph-style workloads directly inside PostgreSQL. The project objectives are:
\begin{itemize}
    \item represent graph data using a clean relational schema,
    \item implement recursive traversal queries using \texttt{WITH RECURSIVE},
    \item compare multiple indexing strategies,
    \item inspect PostgreSQL execution plans using \texttt{EXPLAIN ANALYZE}, and
    \item evaluate how runtime, planner cost, and buffer activity scale with dataset size and query type.
\end{itemize}

\section{Relational Graph Model}
We use two tables:
\begin{itemize}
    \item \texttt{nodes(id, label, node\_type)}
    \item \texttt{edges(src, dst, weight, interaction\_type)}
\end{itemize}

The \texttt{nodes} table stores entities, while \texttt{edges} stores directed relationships between them. This is the adjacency-list view of the graph inside a relational database.

\paragraph{Index configurations.}
We benchmark three index modes:
\begin{enumerate}
    \item \textbf{no\_index}: only primary/foreign-key constraints,
    \item \textbf{basic\_index}: B-tree indexes on \texttt{edges(src)} and \texttt{edges(dst)},
    \item \textbf{composite\_index}: composite B-tree indexes on \texttt{(src, dst)} and \texttt{(dst, src)} in addition to the basic strategy.
\end{enumerate}

\section{Queries Implemented}
The benchmark covers the following workloads:
\begin{itemize}
    \item \textbf{Reachability}: all vertices reachable from a source using recursive expansion.
    \item \textbf{$k$-hop neighborhood}: nodes within bounded depth from a source.
    \item \textbf{Shortest-path BFS style}: bounded recursive path exploration with cycle avoidance.
    \item \textbf{Mutual friends}: self-join based common-neighbor computation.
    \item \textbf{Friend recommendation}: two-hop expansion while excluding direct neighbors.
    \item \textbf{Contact tracing}: bounded recursive propagation query.
\end{itemize}

\section{Experimental Setup}
\subsection{Datasets}
The project uses three synthetic datasets: \texttt{small}, \texttt{medium}, and \texttt{large}. The updated data generator creates graphs with the following structural properties:
\begin{itemize}
    \item a weak ring backbone to avoid complete fragmentation,
    \item community-biased local links,
    \item a small set of hub vertices,
    \item some reciprocal edges and some cross-community edges.
\end{itemize}
This is more informative than a purely uniform random graph because it creates a workload where recursive traversals, self-joins, and recommendation-style queries stress the optimizer differently.

\subsection{Benchmark Methodology}
The updated benchmark script improves over a single-run measurement in four ways:
\begin{enumerate}
    \item each query is tested on multiple parameter cases instead of a single fixed source vertex,
    \item each case uses one warmup run followed by multiple measured runs,
    \item run-level, case-level, and aggregated summaries are all saved, and
    \item plan structure and buffer activity are recorded in addition to execution time.
\end{enumerate}

For every measured query execution, we collect:
\begin{itemize}
    \item planning time,
    \item execution time,
    \item estimated total cost,
    \item shared buffer hits and reads,
    \item plan depth and number of plan nodes,
    \item scan, join, and recursive-union node counts,
    \item representative JSON plan files for inspection.
\end{itemize}

The benchmark produces:
\begin{itemize}
    \item \texttt{results/benchmark\_runs.csv} for detailed run-level results,
    \item \texttt{results/benchmark\_case\_summary.csv} for per-parameter summaries,
    \item \texttt{results/benchmark\_results.csv} for aggregate query-level summaries,
    \item \texttt{results/plans/*.json} for representative execution plans,
    \item \texttt{results/plots/*.png} for visualization.
\end{itemize}

\section{Results}
This section is written to match the updated plot script directly. All referenced figures are generated from the directory structure used by the project.

\subsection{Execution time across query types}
Figures~\ref{fig:exec-small}, \ref{fig:exec-medium}, and \ref{fig:exec-large} compare mean execution time across query families for each dataset. The plots use a log-scale y-axis because recursive traversals and self-joins can differ by more than an order of magnitude.

\begin{figure}[H]
    \centering
    \includegraphics[width=0.95\textwidth]{../results/plots/execution_time_small.png}
    \caption{Mean execution time by query and index mode on the small dataset.}
    \label{fig:exec-small}
\end{figure}

\begin{figure}[H]
    \centering
    \includegraphics[width=0.95\textwidth]{../results/plots/execution_time_medium.png}
    \caption{Mean execution time by query and index mode on the medium dataset.}
    \label{fig:exec-medium}
\end{figure}

\begin{figure}[H]
    \centering
    \includegraphics[width=0.95\textwidth]{../results/plots/execution_time_large.png}
    \caption{Mean execution time by query and index mode on the large dataset.}
    \label{fig:exec-large}
\end{figure}

\paragraph{Discussion.}
These figures typically show the clearest benefit for recursive forward traversals and two-hop recommendation workloads once indexes on \texttt{src} are introduced. Mutual-neighbor queries and recommendation queries may also benefit from better support for repeated self-joins on \texttt{edges}. The large-dataset plot is usually the most informative because it amplifies the difference between scan-heavy and index-supported plans.

\subsection{Speedup relative to the no-index baseline}
Raw runtime plots are useful, but relative speedup communicates the indexing benefit more directly. Figures~\ref{fig:speedup-small}, \ref{fig:speedup-medium}, and \ref{fig:speedup-large} show the speedup of \texttt{basic\_index} and \texttt{composite\_index} relative to \texttt{no\_index}.

\begin{figure}[H]
    \centering
    \includegraphics[width=0.9\textwidth]{../results/plots/speedup_heatmap_small.png}
    \caption{Speedup heatmap on the small dataset. Values larger than 1 indicate improvement over the no-index baseline.}
    \label{fig:speedup-small}
\end{figure}

\begin{figure}[H]
    \centering
    \includegraphics[width=0.9\textwidth]{../results/plots/speedup_heatmap_medium.png}
    \caption{Speedup heatmap on the medium dataset.}
    \label{fig:speedup-medium}
\end{figure}

\begin{figure}[H]
    \centering
    \includegraphics[width=0.9\textwidth]{../results/plots/speedup_heatmap_large.png}
    \caption{Speedup heatmap on the large dataset.}
    \label{fig:speedup-large}
\end{figure}

\paragraph{Discussion.}
These heatmaps make it easy to answer the main design question: \emph{which query classes benefit materially from indexing, and which do not?} In most runs, recursive traversals and self-join workloads should show the largest improvement. If a heatmap entry is close to 1, the index is not helping much for that workload under the chosen data distribution.

\subsection{Scaling behavior}
To study scalability, the updated plotting script generates one plot per query family, with dataset size on the x-axis and runtime on the y-axis. This avoids crowding too many lines in one figure while still making trends visible.

\begin{figure}[H]
    \centering
    \includegraphics[width=0.85\textwidth]{../results/plots/scaling_reachability.png}
    \caption{Scaling behavior for reachability.}
\end{figure}

\begin{figure}[H]
    \centering
    \includegraphics[width=0.85\textwidth]{../results/plots/scaling_k_hop_neighborhood.png}
    \caption{Scaling behavior for bounded neighborhood search.}
\end{figure}

\begin{figure}[H]
    \centering
    \includegraphics[width=0.85\textwidth]{../results/plots/scaling_shortest_path_bfs_style.png}
    \caption{Scaling behavior for BFS-style shortest path search.}
\end{figure}

\begin{figure}[H]
    \centering
    \includegraphics[width=0.85\textwidth]{../results/plots/scaling_mutual_friends.png}
    \caption{Scaling behavior for mutual-friend detection.}
\end{figure}

\begin{figure}[H]
    \centering
    \includegraphics[width=0.85\textwidth]{../results/plots/scaling_friend_recommendation.png}
    \caption{Scaling behavior for friend recommendation.}
\end{figure}

\begin{figure}[H]
    \centering
    \includegraphics[width=0.85\textwidth]{../results/plots/scaling_contact_tracing.png}
    \caption{Scaling behavior for contact tracing.}
\end{figure}

\paragraph{Discussion.}
These plots reveal whether indexing merely reduces the constant factor or also improves how gracefully runtime grows with dataset size. Recursive queries usually remain the most expensive as the graph grows, but indexes can significantly reduce the slope of the growth curve.

\subsection{Planner cost, measured runtime, and buffer activity}
Runtime alone does not explain \emph{why} one configuration is better. The updated benchmark therefore also records planner cost and buffer behavior.

\begin{figure}[H]
    \centering
    \includegraphics[width=0.82\textwidth]{../results/plots/cost_vs_runtime_all.png}
    \caption{Relationship between PostgreSQL's estimated total cost and measured runtime.}
\end{figure}

\begin{figure}[H]
    \centering
    \includegraphics[width=0.9\textwidth]{../results/plots/shared_reads_large.png}
    \caption{Mean shared read blocks by query and index mode on the large dataset.}
\end{figure}

\begin{figure}[H]
    \centering
    \includegraphics[width=0.9\textwidth]{../results/plots/shared_hits_large.png}
    \caption{Mean shared hit blocks by query and index mode on the large dataset.}
\end{figure}

\paragraph{Discussion.}
The cost-vs-runtime scatter plot helps us check whether PostgreSQL's optimizer cost is directionally aligned with measured runtime. The buffer plots show whether an index mode reduces expensive block reads, improves cache locality, or both. In many workloads, a lower runtime under indexing coincides with a substantial reduction in shared read blocks.

\subsection{Planning overhead and parameter sensitivity}
Two additional analyses are useful for recursive SQL workloads:
\begin{itemize}
    \item \textbf{Planning-to-execution ratio}, which reveals whether planner overhead is negligible or significant compared to actual execution.
    \item \textbf{Case-wise variation}, which captures how much the same query family changes with different source/destination choices or recursion depth.
\end{itemize}

\begin{figure}[H]
    \centering
    \includegraphics[width=0.9\textwidth]{../results/plots/planning_ratio_by_query.png}
    \caption{Average planning-to-execution ratio by query family and index mode.}
\end{figure}

\begin{figure}[H]
    \centering
    \includegraphics[width=0.9\textwidth]{../results/plots/case_variation_k_hop_neighborhood_large.png}
    \caption{Case-wise runtime variation for bounded neighborhood search on the large dataset.}
\end{figure}

\begin{figure}[H]
    \centering
    \includegraphics[width=0.9\textwidth]{../results/plots/case_variation_shortest_path_bfs_style_large.png}
    \caption{Case-wise runtime variation for BFS-style shortest path on the large dataset.}
\end{figure}

\begin{figure}[H]
    \centering
    \includegraphics[width=0.9\textwidth]{../results/plots/case_variation_contact_tracing_large.png}
    \caption{Case-wise runtime variation for contact tracing on the large dataset.}
\end{figure}

\paragraph{Discussion.}
A single source vertex can be misleading in graph workloads. These case-wise plots justify the updated methodology: performance can vary significantly across parameter settings, especially for bounded recursive search. Averaging across several representative cases gives a more reliable performance picture.

\section{Use Cases}
\subsection{Social Network Analysis}
In a social network, users are nodes and interactions are edges. Reachability can answer multi-hop connectivity questions, mutual-friend queries support common-neighbor discovery, and friend recommendation naturally maps to two-hop traversal with direct-neighbor filtering.

\subsection{Disease Spread and Contact Tracing}
In a contact graph, individuals are nodes and physical interactions are edges. Bounded recursive traversal models contact tracing by identifying all individuals within a small number of hops from an infected source. Because such queries repeatedly expand frontiers, they are a good fit for recursive CTEs and an informative stress test for the optimizer.

\section{Main Findings}
The final conclusions should be supported directly by the generated CSV files and plots. In most runs, the following patterns are expected:
\begin{itemize}
    \item basic indexes already provide a large improvement for forward traversal queries,
    \item composite indexes can further help join-heavy workloads, though the improvement may not be uniform,
    \item recursive queries become increasingly expensive with dataset size and recursion depth,
    \item buffer-read reductions often explain a significant part of the observed speedup,
    \item planner cost is usually correlated with runtime, but not perfectly.
\end{itemize}

\paragraph{Optional quantitative summary.}
Before submission, add one short table using \texttt{results/best\_index\_mode\_by\_query.csv} to report the fastest index mode for each query on each dataset. This makes the conclusion section more concrete without requiring a large amount of manual tabulation.

\section{Conclusion}
This project demonstrates that graph-style workloads can be expressed and executed effectively in PostgreSQL using recursive SQL and careful indexing. A relational system is not a drop-in replacement for a dedicated graph engine for every possible workload, but it is a practical and often efficient solution for moderate graph analytics when schema design, traversal formulation, and indexing strategy are chosen carefully.

\end{document}
